\begin{theorem}[Proved: exponent-two normal nodes with arbitrary dyadic ramification]
\label{C:ramified_exponent_two_node}
Let $A$ be any characteristic-zero DVR of residue characteristic two,
let $\pi$ be a uniformizer, and let $u\in A^\times$. Put
\[
 R=\left(A[x,y]/(xy-\pi^2u)\right)_{(\pi,x,y)}
 \quad\hbox{or}\quad
 R=A[[x,y]]/(xy-\pi^2u).                              \tag{C.1.RN1}
\]
Both rings are flat normal two-dimensional singular local domains.
For either ring the degree-two polynomial map is not surjective,
and the generic comparison has a nonzero kernel:
\[
 \ker\bigl(R_R(C_2)\longrightarrow R_{R[1/2]}(C_2)\bigr)\ne0.
                                                               \tag{C.1.RN2}
\]
If the ramification index $e=v_\pi(2)$ is greater than one,
a rank-two involution gives an explicit class in this kernel.
For the polynomial localization, $K_0$ of every nonzero integer
localization is $\mathbb Z$. No generic $K_0$ assertion for the
formal ring is required.
\end{theorem}
\begin{proof}
The ring-theoretic arguments in
Theorem~\ref{C:dyadic_normal_node_criterion} apply with the
uniformizer $\pi$ replacing two. We spell out the hypotheses
to ensure that no completeness of $A$ is inserted.
The polynomial equation is primitive and irreducible over
$A[x]$, so its localization is a two-dimensional domain,
flat over the DVR. It is a hypersurface and hence satisfies
$(S_2)$. Its only vertical height-one primes are $(\pi,x)$
and $(\pi,y)$; their localizations have maximal ideal $(\pi)$
and are DVRs. Away from $\pi$ the ring is a localization of
$\operatorname{Frac}(A)[x,x^{-1}]$, a PID. Thus it satisfies
$(R_1)$ and is normal by Serre's criterion
\cite[Lemma 10.157.4, Tag 031S]{C-NormalKernelStacks}.

For the formal ring, $S=A[[x,y]]$ is a regular local ring
of dimension three, hence a UFD
\cite[Lemma 15.123.2, Tag 0AG0]{C-NormalNodesStacks}.
If $xy-\pi^2u=gh$ with both factors nonunits, reduction
modulo $\pi$ forces, after interchanging them,
$\bar g=xv$ and $\bar h=yw$ with $v,w$ units.
Setting $x=0$ contradicts unique factorization of
$g(0,y)h(0,y)=-\pi^2u$ in $A[[y]]$: $h(0,y)$ is
neither a unit nor divisible by $\pi$.
Thus the quotient is a domain of dimension two and is flat
over $A$. It is a hypersurface, and its two vertical
height-one localizations are DVRs as above.
At every prime away from $\pi$, $x$ is a unit and the
derivation $\partial/\partial y$ sends the defining equation
to $x$. The equation therefore lies outside the square of
the maximal ideal of the corresponding regular localization
of $S$, so that quotient is regular. This proves $(R_1)$,
and hence normality, also in the formal case.
In both cases the relation lies in the square of the
ambient maximal ideal, so the embedding dimension is three
and the dimension is two; the ring is singular.
For the polynomial ring, locality and the displayed generic
PID give the $K_0$ assertion, since odd integers are already
units and every even integer localization inverts $\pi$.

Write $2=v\pi^e$ with $v\in A^\times$.
If $e=1$, the equation is $xy=4u/v^2$ and the entire
assertion about the kernel and polynomial map is already
Theorem~\ref{C:dyadic_normal_node_criterion}, with its
explicit rank-six witness. Assume now that $e>1$ and put
\[
 h=1+2/\pi=1+v\pi^{e-1}\in A^\times,
 \qquad a=1+\pi,\qquad
 T=\begin{pmatrix}a&x\\-h u^{-1}y&-a\end{pmatrix}.
                                                               \tag{C.1.RN3}
\]
The identity $a^2-1=\pi^2h$ and the node equation give
$T^2=I_2$. We will prove that
$\eta=[T]-\mathbf1-\mathrm{sgn}$ is nonzero and lies
outside the degree-two polynomial image.

Let $\widehat A$ denote the completion of $A$ and put
$K=\widehat A[\epsilon]/(\epsilon^2)$.
There is a ring map
\[
 R\longrightarrow K,\qquad
 x\longmapsto\pi+\epsilon,\qquad
 y\longmapsto\pi u-u\epsilon.                         \tag{C.1.RN4}
\]
The product of the two images is $\pi^2u$.
For the polynomial localization, every inverted denominator
has nonzero residue in $A/\pi A$ and is therefore a unit in
the local target. For the formal ring this is a convergent
evaluation: $K$ is complete for $(\pi,\epsilon)$, the images
of both variables belong to that ideal, and its topology
agrees with the $\pi$-adic topology since $\epsilon^2=0$.
In degree $i+j$, the constant term is divisible by
$\pi^{i+j}$ and the $\epsilon$ coefficient by
$\pi^{i+j-1}$, so the defining infinite sums converge.
Thus (C.1.RN4) is valid even when $A$ is not complete.

Under this map the involution becomes $T_0+\epsilon U$, where
\[
 T_0=\begin{pmatrix}1+\pi&\pi\\-h\pi&-1-\pi\end{pmatrix},
 \qquad U=\begin{pmatrix}0&1\\h&0\end{pmatrix}.         \tag{C.1.RN5}
\]
These matrices satisfy
$T_0^2=I$, $T_0U+UT_0=0$ and $U^2=hI$.
Set $w_+=(1,-1)^t$ and $w_-=(1,-h)^t$.
The difference and norm maps factor as
\[
 T_0-I=\pi w_-\begin{pmatrix}1&1\end{pmatrix},
 \qquad
 T_0+I=\pi w_+\begin{pmatrix}h&1\end{pmatrix}.
                                                               \tag{C.1.RN6}
\]
Consequently the plus cycles are $\widehat A w_+$ and
the minus cycles are $\widehat A w_-$, while the respective
norm and difference images are their multiples by $\pi$.
Both Tate groups are therefore the one-dimensional residue
field $F=A/\pi A$.
Moreover
\[
 Uw_+=-w_-,\qquad Uw_-=-h w_+.
\]
For the coefficient functional sending $\epsilon$ to $1$,
the two odd linear-nil Bockstein maps therefore have square
equal to multiplication by $\bar h=1$ on these Tate groups.

Apply the finite-length primary parity invariant proved in
Theorem~\ref{C:ramified_squarezero_criterion}, which extends
Lemma~\ref{C:linear_nil_Bockstein} to every dyadic DVR.
For the nonzero primary $q(z)=z-1$ its even Tate primary
has composition length one. Hence
\[
             \chi_{z-1,\ell}(T_0+\epsilon U)=1,
 \qquad \chi_{z-1,\ell}(\mathbf1)
       =\chi_{z-1,\ell}(\mathrm{sgn})=0.                \tag{C.1.RN7}
\]
The constant lines have zero odd Bockstein map, so their
$z-1$ primary is zero, including when their full Tate
groups have length greater than one.
The invariant is additive on every admissible exact sequence,
by the exact Tate cycle and primary decomposition proved
in that cited theorem. Base change along (C.1.RN4) preserves
such sequences because their underlying projective-module
sequences split. Thus (C.1.RN7) detects the source class
$\eta$ without requiring the map (C.1.RN4) to be flat.

The local ring $R$ has only free finite projective modules.
Its degree-two polynomial image is the span of the trivial
and sign lines, using the admissible regular-module sequence.
The displayed detector vanishes on that span and takes value
one on $[T]$, proving failure of degree-two surjectivity.

Finally, after inverting two, $\pi$ and $x$ are units, and
\[
 P=\begin{pmatrix}x&x\\1-a&-1-a\end{pmatrix},
 \qquad \det P=-2x,\qquad
 TP=P\operatorname{diag}(1,-1).
\]
Both generic eigenlines are explicitly free. Therefore
$\eta$ lies in the kernel (C.1.RN2), as claimed.
No order of this source class is asserted.
\end{proof}

This theorem settles exponent two for every dyadic ramification
index, together with the preceding unramified whole-family
criterion. It does not assert an all-exponent criterion over
ramified DVRs.

The exact certificate \texttt{code/C\_finish\_ramified\_nodes.py}
checks twelve symbolic identity groups and forty-two square-zero
specializations over
$\mathbf Z_2[\pi]/(\pi^e-2,4)$ for $2\le e\le8$ and explicit
residue-unit choices. Its output and runtime are saved in
\texttt{logs/C\_finish\_ramified\_nodes.json}.
These are bounded matrix and residue checks. Normality, convergence
of formal evaluation and arbitrary-rank additivity of the detector
are established by the written proof.
